\section{AMC 12 12A 2014}
        \begin{enumerate}
        	\item (\textit{AMC 12 12A 2014}) What is $10 \cdot \left(\frac{1}{2} + \frac{1}{5} + \frac{1}{10}\right)^{-1}?$
$\textbf{(A)}\ 3\qquad\textbf{(B)}\ 8\qquad\textbf{(C)}\ \frac{25}{2}\qquad\textbf{(D)}\ \frac{170}{3}\qquad\textbf{(E)}\ 170$

 \textbf{Answer: }C



\textbf{Solution}

We have 
 \[10\cdot\left(\frac{1}{2}+\frac{1}{5}+\frac{1}{10}\right)^{-1}\]Making the denominators equal gives
 \[\implies 10\cdot\left(\frac{5}{10}+\frac{2}{10}+\frac{1}{10}\right)^{-1}\]
 \[\implies 10\cdot\left(\frac{5+2+1}{10}\right)^{-1}\]
 \[\implies 10\cdot\left(\frac{8}{10}\right)^{-1}\]
 \[\implies 10\cdot\left(\frac{4}{5}\right)^{-1}\]
 \[\implies 10\cdot\frac{5}{4}\]
 \[\implies \frac{50}{4}\]Finally, simplifying gives
 \[\implies \boxed{\textbf{(C)}\ \frac{25}{2}}\]




\textbf{Solution 2}

We have
 \[\left(\frac{1}{10}\right)^{-1}\cdot \left(\frac{1}{2} + \frac{1}{5} + \frac{1}{10}\right)^{-1}\]By Distributive Property,
 \[\left(\frac{1}{20}+\frac{1}{50}+\frac{1}{100}\right)^{-1}\]Now, we want to find the least common multiple of $20, 50,$ and $100,$ so
 \[\text{lcm}(20,50,100)=\text{lcm}(2^2 \cdot 5,2 \cdot 5^2,2^2 \cdot 5^2)=2^2 \cdot 5^2=100\]Converting everything to a denominator of $100,$
 \[\left(\frac{5}{100}+\frac{2}{100}+\frac{1}{100}\right)^{-1}=\left(\frac{8}{100}\right)^{-1}=\frac{100}{8}\]Now, we use Euclidean Algorithm, to find if this fraction is reducible, so
 \[\gcd(100,8)=\gcd(12,8)=\gcd(4,8)=\gcd(4,4)\]Thus, both the numerator and denominator are divisible by $4,$ so
 \[\frac{100}{8} \cdot \frac{4}{4}=\frac{100}{4} \cdot \frac{4}{8}=25 \cdot \frac{1}{2}=\boxed{\frac{25}{2}}\]

 - kante314

 


	\item (\textit{AMC 12 12A 2014}) At the theater children get in for half price.  The price for $5$ adult tickets and $4$ child tickets is $\$24.50$.  How much would $8$ adult tickets and $6$ child tickets cost?


 $\textbf{(A) }\$35\qquad \textbf{(B) }\$38.50\qquad \textbf{(C) }\$40\qquad \textbf{(D) }\$42\qquad \textbf{(E) }\$42.50$

 \textbf{Answer: }B



\textbf{Solution}

Suppose $x$ is the price of an adult ticket. The price of a child ticket would be $\frac{x}{2}$.

 
 \begin{eqnarray*} 5x + 4(x/2) = 7x &=& 24.50\\ x &=& 3.50\\ \end{eqnarray*}

 Plug in for 8 adult tickets and 6 child tickets.

 
 \begin{eqnarray*} 8x + 6(x/2) &=& 8(3.50) + 3(3.50)\\  &=&\boxed{\textbf{(B)}\ \ 38.50}\\ \end{eqnarray*}



	\item (\textit{AMC 12 12A 2014}) Walking down Jane Street, Ralph passed four houses in a row, each painted a different color. He passed the orange house before the red house, and he passed the blue house before the yellow house. The blue house was not next to the yellow house. How many orderings of the colored houses are possible?


 $\textbf{(A)}\ 2\qquad\textbf{(B)}\ 3\qquad\textbf{(C)}\ 4\qquad\textbf{(D)}\ 5\qquad\textbf{(E)}\ 6$

 \textbf{Answer: }B



\textbf{Solution 1}

Let's use casework on the yellow house. The yellow house $(\text{Y})$ is either the $3^\text{rd}$ house or the last house.

 Case 1: $\text{Y}$ is the $3^\text{rd}$ house.

 The only possible arrangement is $\text{B}-\text{O}-\text{Y}-\text{R}$

 Case 2: $\text{Y}$ is the last house.

 There are two possible arrangements:

 $\text{B}-\text{O}-\text{R}-\text{Y}$
$\text{O}-\text{B}-\text{R}-\text{Y}$ 

 The answer is $1+2=\boxed{\textbf{(B) } 3}$




\textbf{Solution 2 (symmetry)}

There are $4!=24$ arrangements without restrictions. There are $3!\cdot2!=12$ arrangements such that the blue house neighboring the yellow house (calculating the arrangments of [$\text{BY}$], $\text{O}$, and $\text{R}$). Hence, there are $24-12=12$ arrangements with the blue and yellow houses non-adjacent.

 By symmetry, exactly half of the $12$ arrangements have the blue house before the yellow house, and exactly half of those $6$ arrangements have the orange house before the red house, so our answer is $12\cdot\frac{1}{2}\cdot\frac{1}{2}= \boxed{\textbf{(B) } 3}$




\textbf{Solution 3}

To start with, the blue house is either the first or second house.

 If the blue house is the first, then the orange must follow, leading to $2$ cases: $\text{B-O-R-Y}$ and $\text{B-O-Y-R}$. 

 If the blue house is second, then the orange house must be first and the yellow house last, leading to $1$ case: $\text{O-B-R-Y}$.Therefore, our answer is $\boxed{\textbf{(B) } 3}$.

 



\textbf{Solution 4 (Complementary Counting)}

We first count all the cases with the restrictions that B comes before Y and O comes before R. We have $4$ "slots" to choose from, and when we choose $2$ to be B and Y, the order for B and Y is automatically chosen. Also, O and R are also automatically chosen, so there are simply $\dbinom42=6$ total cases. 

 From here, the cases that don't work are the ones where B and Y are adjacent. BY can be treated as a single block; now, we have $3$ slots and $1$ thing to place, so there are $3$ of these cases. 

 Thus, this yields $6-3=3$ total cases. We are done. 

 


	\item (\textit{AMC 12 12A 2014}) Suppose that $a$ cows give $b$ gallons of milk in $c$ days. At this rate, how many gallons of milk will $d$ cows give in $e$ days?


 $\textbf{(A)}\ \frac{bde}{ac}\qquad\textbf{(B)}\ \frac{ac}{bde}\qquad\textbf{(C)}\ \frac{abde}{c}\qquad\textbf{(D)}\ \frac{bcde}{a}\qquad\textbf{(E)}\ \frac{abc}{de}$

 \textbf{Answer: }A



\textbf{Solution 1}

We need to multiply $b$ by $\frac{d}{a}$ for the new cows and $\frac{e}{c}$ for the new time, so the answer is $b\cdot \frac{d}{a}\cdot \frac{e}{c}=\frac{bde}{ac}$, or $\boxed{\textbf{(A) } \frac{bde}{ac}}$.

 



\textbf{Solution 2}

We plug in $a=2$, $b=3$, $c=4$, $d=5$, and $e=6$. Hence the question becomes "2 cows give 3 gallons of milk in 4 days. How many gallons of milk do 5 cows give in 6 days?"

 If 2 cows give 3 gallons of milk in 4 days, then 2 cows give $\frac{3}{4}$ gallons of milk in 1 day, so 1 cow gives $\frac{3}{4\cdot2}$ gallons in 1 day. This means that 5 cows give $\frac{5\cdot3}{4\cdot2}$ gallons of milk in 1 day. Finally, we see that 5 cows give $\frac{5\cdot3\cdot6}{4\cdot2}$ gallons of milk in 6 days. Substituting our values for the variables, this becomes $\frac{dbe}{ac}$, which is $\boxed{\textbf{(A)}\ \frac{bde}{ac}}$.

 



\textbf{Solution 3}

We see that the number of cows is inversely proportional to the number of days and directly proportional to the gallons of milk. So our constant is $\dfrac{ac}{b}$.

 Let $g$ be the answer to the question. We have $\dfrac{de}{g}=\dfrac{ac}{b}\implies gac=bde\implies g=\dfrac{bde}{ac}\implies\boxed{ \textbf{(A)}\ \frac{bde}{ac}}$




\textbf{Solution 4}

The problem specifics "rate," so it would be wise to first find the rate at which cows produce milk. We can find the rate of work/production by finding the gallons produced by a single cow in a single day. To do this, we divide the amount produced by the number of cows and number of days 
 \[\implies\text{rate}=\dfrac{b}{ac}\]

 Now that we have the gallons produced by a single cow in a single day, we simply multiply that rate by the number of cows and the number of days 
 \[\implies\boxed{\textbf{(A) } \dfrac{bde}{ac}}\]




\textbf{Solution 5}

If $a$ cows give $b$ gallons of milk in $c$ days, that means that one cow will give $\frac{b}{a}$ gallons of milk in $c$ days. Also, we want to find the number of gallons of milk will $d$ cows give in $e$ days, so in $\frac{e}{c}$ days $d$ cows give $\frac{bd}{a}$ gallons of milk. Multiplying with the formula $d=rt$, we get $\boxed{\textbf{(A) } \dfrac{bde}{ac}}$




\textbf{Solution 6}

In the right formula, plugging in $d=a$ and $e=c$ should simplify to $b$, as if it doesn't, we'd essentially be saying "$a$ cows give $b$ gallons of milk in $c$ days, but $a$ cows don't give $b$ gallons of milk in $c$ days." The only one of the answer choices that simplifies like this is $\boxed{\textbf{(A)}}$




\textbf{Solution 7}

Call the rate at which a single cow produces milk $r$. Using the given information we get that $rac=b$, and we want to find $rde$. We solve for $r$ in the first equation to get $r=b \cdot \frac{1}{ac}$. Plugging in this value for $r$ into the expression we want, we get that $rde=\frac{bde}{ac}$, or $\boxed{\textbf{(A)}}$.'

 



\textbf{Solution 8 (Formula Version)}

We can start by using the work formula $W = n \cdot t \cdot r$, where $W$ is the total work done (milk produced), $n$ is the number of workers (cows), $t$ is the time (days), and $r$ is the rate of a single worker. (Note: I learned this special formula from the Aops book "The Basics" by Richard Rusczyk and I find it very useful and convenient when solving these problems)

 From the first part of the problem, we know that $a$ cows produce $b$ gallons of milk in $c$ days. Substituting $n = a$, $W = b$, and $t = c$ into our formula gives:$b = a \cdot c \cdot r \implies r = \frac{b}{ac}$

 And just like that, we can continue by using the ways from the previous solutions to find the answer $\boxed{\textbf{(A)}}$.

 


	\item (\textit{AMC 12 12A 2014}) On an algebra quiz, $10\%$ of the students scored $70$ points, $35\%$ scored $80$ points, $30\%$ scored $90$ points, and the rest scored $100$ points. What is the difference between the mean and median score of the students' scores on this quiz?


 $\textbf{(A)}\ 1\qquad\textbf{(B)}\ 2\qquad\textbf{(C)}\ 3\qquad\textbf{(D)}\ 4\qquad\textbf{(E)}\ 5$

 \textbf{Answer: }C



\textbf{Solution 1}

WLOG, let there be $20$ students (the least whole number possible) who took the test. We have $2$ students score $70$ points, $7$ students score $80$ points, $6$ students score $90$ points and $5$ students score $100$ points. Therefore, the meanis $87$ and the median is $90$.

 Thus, the solution is 
 \[90-87=3\implies\boxed{\textbf{(C)} \ 3}\]




\textbf{Solution 2}

The percentage who scored $100$ points is $100\%-(10\%+35\%+30\%)=100\%-75\%=25\%$. Now, we need to find the median, which is the score that splits the upper and lower $50\%$.The lower $10\%+35\%=45\%$ scored $70$ or $80$ points, so the median is $90$ (since the upper $25\%$ is $100$ points and the lower $45\%$ is $70$ or $80$).The mean is $10\%\cdot70+35\%\cdot80+30\%\cdot90+25\%\cdot100=7+28+27+25=87$.So, our solution is $90-87=3\Rightarrow\boxed{3 \ \textbf{(C)} }$ ~sosiaops

 



\textbf{Solution 3}

The $\le 80$-point scores make up $10\%+35\% = 45\% < 50\%$ of the scores, but the $\le 90$-point scores make up $45\%+30\% = 75\% > 50\%$ of the scores, so the median is $90$.

 $10\%$ of scores were $70-90 = -20$ more than the median, $35\%$ were $-10$ more, $100\%-75\% = 25\%$ were $10$ more, and the rest were equal. This means that the mean score is $10\%\cdot(-20)+35\%\cdot(-10)+25\%\cdot10 = -2 + (-3.5) + 2.5 = -3$ more than the median, so their difference is $\left|-3\right| = \boxed{\textbf{(C)}\ 3}$.~\href{https://artofproblemsolving.com/wiki/index.php?title=User:Emerald\_block}{emerald\_block}

 


	\item (\textit{AMC 12 12A 2014}) The difference between a two-digit number and the number obtained by reversing its digits is $5$ times the sum of the digits of either number.  What is the sum of the two digit number and its reverse?


 $\textbf{(A) }44\qquad \textbf{(B) }55\qquad \textbf{(C) }77\qquad \textbf{(D) }99\qquad \textbf{(E) }110$

 \textbf{Answer: }D



\textbf{Solution 1}

Let the two digits be $a$ and $b$. Then, $5a + 5b = 10a + b - 10b - a = 9a - 9b$, or $2a = 7b$. This yields $a = 7$ and $b = 2$ because $a, b < 10$. Then, $72 + 27 = \boxed{\textbf{(D) }99}.$




\textbf{Solution 2 (Meta)}

We start like above. Let the two digits be $a$ and $b$. Therefore, $5(a+b) = 10a+b-10b-a=9(a-b)$. Since we are looking for $10a+b+10b+a=11(a+b)$ and we know that $a+b$ must be a multiple of $9$, the only answer choice that works is $\boxed{\textbf{(D) }99}.$



	\item (\textit{AMC 12 12A 2014}) The first three terms of a geometric progression are $\sqrt 3$, $\sqrt[3]3$, and $\sqrt[6]3$.  What is the fourth term?


 $\textbf{(A) }1\qquad \textbf{(B) }\sqrt[7]3\qquad \textbf{(C) }\sqrt[8]3\qquad \textbf{(D) }\sqrt[9]3\qquad \textbf{(E) }\sqrt[10]3\qquad$

 \textbf{Answer: }A



\textbf{Solution}

The terms are  $\sqrt 3$, $\sqrt[3]3$, and $\sqrt[6]3$, which are equivalent to $3^{\frac{3}{6}}$, $3^{\frac{2}{6}}$, and $3^{\frac{1}{6}}$.  So the next term will be $3^{\frac{0}{6}}=1$, so the answer is $\boxed{\textbf{(A)}}$.

 


	\item (\textit{AMC 12 12A 2014}) A customer who intends to purchase an appliance has three coupons, only one of which may be used:


 Coupon 1: $10\%$ off the listed price if the listed price is at least $\$50$


 Coupon 2: $\$20$ off the listed price if the listed price is at least $\$100$


 Coupon 3: $18\%$ off the amount by which the listed price exceeds $\$100$


 For which of the following listed prices will coupon $1$ offer a greater price reduction than either coupon $2$ or coupon $3$?


 $\textbf{(A) }\$179.95\qquad \textbf{(B) }\$199.95\qquad \textbf{(C) }\$219.95\qquad \textbf{(D) }\$239.95\qquad \textbf{(E) }\$259.95\qquad$

 \textbf{Answer: }C



\textbf{Solution 1}

Let the listed price be $x$. Since all the answer choices are above $\textdollar100$, we can assume $x > 100$. Thus the discounts after the coupons are used will be as follows:

 Coupon 1: $x\times10\%=.1x$

 Coupon 2: $20$

 Coupon 3: $18\%\times(x-100)=.18x-18$

 
For coupon $1$ to give a greater price reduction than the other coupons, we must have $.1x>20\implies x>200$ and $.1x>.18x-18\implies.08x<18\implies x<225$.

 The only choice that satisfies such conditions is $\boxed{\textbf{(C)}\ \textdollar219.95}$




\textbf{Solution 2 (Using The Answers)}

For coupon $1$ to be the most effective, we want 10\% of the price to be greater than 20. This clearly occurs if the value is over 200. For coupon 1 to be more effective than coupon 3, we want to minimize the value over 200, so $\boxed{\textbf{(C) }\textdollar219.95}$ is the smallest answer choice over 200.

 



\textbf{Solution 3}

Since all the answer choices are above $\$100$, let the price be $\$100 + x$, where $x > 0$. Then, the discount from Coupon 1 is $\$10 + \frac{x}{10}$, the discount from Coupon 2 is $\$20$, and the discount from Coupon 3 is $\frac{9}{50}x$. Thus we obtain the following inequalities: \begin{align*}&10 + \frac{x}{10} > 20 \\&10 + \frac{x}{10} > \frac{9}{50}x\end{align*}
Solving them, we obtain $x > 100$ and $x < 125$. Thus the price is between $\$200$ and $\$225\Rightarrow \boxed{\text{(C)}}$.

 


	\item (\textit{AMC 12 12A 2014}) Five positive consecutive integers starting with $a$ have average $b$. What is the average of $5$ consecutive integers that start with $b$?


 $\textbf{(A)}\ a+3\qquad\textbf{(B)}\ a+4\qquad\textbf{(C)}\ a+5\qquad\textbf{(D)}\ a+6\qquad\textbf{(E)}\ a+7$

 \textbf{Answer: }B



\textbf{Solution 1}

Let $a=1$. Our list is $\{1,2,3,4,5\}$ with an average of $15\div 5=3$. Our next set starting with $3$ is $\{3,4,5,6,7\}$. Our average is $25\div 5=5$.

 Therefore, we notice that $5=1+4$ which means that the answer is $\boxed{\textbf{(B)}\ a+4}$.

 



\textbf{Solution 2}

We are given that 
 \[b=\frac{a+a+1+a+2+a+3+a+4}{5}\]
 \[\implies b =a+2\]

 We are asked to find the average of the 5 consecutive integers starting from $b$ in terms of $a$. By substitution, this is 
 \[\frac{a+2+a+3+a+4+a+5+a+6}5=a+4\]

 Thus, the answer is $\boxed{\textbf{(B)}\ a+4}$




\textbf{Solution 3 (fast)}

We know from experience (or logical reasoning) that the average of $5$ consecutive numbers is the $3^\text{rd}$ one or the $1^\text{st} + 2$. With the logic, we find that $b=a+2$. Therefore, $b+2=(a+2)+2=\boxed{a+4}$.

 



\textbf{Solution 4}

The list of numbers is $\left\{a,\ a+1,\ b,\ a+3,\ a+4\right\}$ so $b=a+2$. The new list is $\left\{a+2,\ a+3,\ a+4,\ a+5,\ a+6\right\}$ and the average is $a+4 \Longrightarrow \boxed{\textbf{(B) } a+4}$.

 


	\item (\textit{AMC 12 12A 2014}) Three congruent isosceles triangles are constructed with their bases on the sides of an equilateral triangle of side length $1$.  The sum of the areas of the three isosceles triangles is the same as the area of the equilateral triangle.  What is the length of one of the two congruent sides of one of the isosceles triangles?


 $\textbf{(A) }\dfrac{\sqrt3}4\qquad \textbf{(B) }\dfrac{\sqrt3}3\qquad \textbf{(C) }\dfrac23\qquad \textbf{(D) }\dfrac{\sqrt2}2\qquad \textbf{(E) }\dfrac{\sqrt3}2$

 \textbf{Answer: }B



\textbf{Solution 1}

Reflect each of the triangles over its respective side.  Then since the areas of the triangles total to the area of the equilateral triangle, it can be seen that the triangles fill up the equilateral one and the vertices of these triangles concur at the circumcenter of the equilateral triangle.  Hence the desired answer is just its circumradius, or $\boxed{\dfrac{\sqrt3}3\textbf{ (B)}}$.

 (Solution by djmathman)

 



\textbf{Solution 2}

Since the total area of each \href{https://artofproblemsolving.com/wiki/index.php?title=Congruent\_(geometry)}{congruent} isosceles triangle is the same, the area of each is $\dfrac{1}3$ the total \href{https://artofproblemsolving.com/wiki/index.php?title=Area\_of\_an\_equilateral\_triangle}{area of the equilateral triangle} of side length 1, or $\dfrac{1}3$ x $\dfrac{\sqrt3}4$. Likewise, the area of each can be defined as $\dfrac{bh}2$ with base $b$ equaling 1, meaning that $\dfrac{h}2$ = $\dfrac{1}3$ x $\dfrac{\sqrt3}4$, or $h$ = $\dfrac{\sqrt3}6$. A side length of the isosceles triangle is the hypotenuse with legs $\dfrac{b}2$ and $h$. Using the Pythagorean Theorem, the side length is $\sqrt{{(\dfrac{1}2)}^2 + {(\dfrac{\sqrt3}6)}^2}$, or $\boxed{\dfrac{\sqrt3}3\textbf{ (B)}}$.

 (Solution by johnstucky)

 



\textbf{Solution 3}

Note that each of the congruent sides meets at the circumcenter of the equilateral triangle. Therefore, they are equal to the circumradius. The area is then $\dfrac{abc}{4R}$, or $\dfrac{1}{4R}$. The area can also be found using $\dfrac{\sqrt3}{4}s^2$, or $\dfrac{\sqrt3}{4}$. Setting these equal and solving entails $R$ = $\dfrac{1}{\sqrt3}$ or $\boxed{\dfrac{\sqrt3}3\textbf{ (B)}}$.

 solution by Proper

 


 
Additional Insight:After noting that the congruent sides meet at the circumcenter, you can just use the Extended Law of Sines to see that $R = \dfrac{1}{2\sin 60^{\circ}} = \dfrac{\sqrt{3}}{3}$, which is choice B~Extremelysupercooldude

 


	\item (\textit{AMC 12 12A 2014}) David drives from his home to the airport to catch a flight.  He drives $35$ miles in the first hour, but realizes that he will be $1$ hour late if he continues at this speed.  He increases his speed by $15$ miles per hour for the rest of the way to the airport and arrives $30$ minutes early.  How many miles is the airport from his home?


 $\textbf{(A) }140\qquad \textbf{(B) }175\qquad \textbf{(C) }210\qquad \textbf{(D) }245\qquad \textbf{(E) }280\qquad$

 \textbf{Answer: }C



\textbf{Solution 1}

Note that he drives at $50$ miles per hour after the first hour and continues doing so until he arrives.

 Let $d$ be the distance still needed to travel after $1$ hour. We have that $\dfrac{d}{50}+1.5=\dfrac{d}{35}$, where the $1.5$ comes from $1$ hour late decreased to $0.5$ hours early.

 Simplifying gives $7d+525=10d$, or $d=175$.

 Now, we must add an extra $35$ miles traveled in the first hour, giving a total of $\boxed{\textbf{(C) } 210}$ miles.

 



\textbf{Solution 2 (Answer Choices)}

Instead of spending time thinking about how one can set up an equation to solve the problem, one can simply start checking the answer choices.Quickly checking, we know that neither choice $\textbf{(A)}$ or choice $\textbf{(B)}$ work, but $\textbf{(C)}$ does. We can verify as follows. After $1$ hour at $35 \text{ mph}$, David has $175$ miles left. This then takes him $3.5$ hours at $50 \text{ mph}$. But $210/35 = 6 \text{ hours}$. Since $1+3.5 = 4.5 \text{ hours}$ is $1.5$ hours less than $6$, our answer is $\boxed{\textbf{(C)} \: 210}$.

 



\textbf{Solution 3}

Let the total distance be $d$, and let the time taken after the speed change be $t$. Then $d=35(t+1)$. Since 1 hour has passed, and he increased his speed by $15$ miles per hour, and he had already traveled $35$ miles, the new equation is $d-35=50(t-1-\frac{1}{2})=50(t-\frac{3}{2})$. Solving, $d=35t+35=50t-40$, $15t=75$, $t=5$, and $d=35(5+1)=35\cdot 6=210 \Longrightarrow \boxed{\textbf{(C) } 210}$.

 


	\item (\textit{AMC 12 12A 2014}) Two circles intersect at points $A$ and $B$.  The minor arcs $AB$ measure $30^\circ$ on one circle and $60^\circ$ on the other circle.  What is the ratio of the area of the larger circle to the area of the smaller circle?


 $\textbf{(A) }2\qquad \textbf{(B) }1+\sqrt3\qquad \textbf{(C) }3\qquad \textbf{(D) }2+\sqrt3\qquad \textbf{(E) }4\qquad$

 \textbf{Answer: }D



\textbf{Solution 1}

Let the radius of the larger and smaller circles be $x$ and $y$, respectively. Also, let their centers be $O_1$ and $O_2$, respectively. Then the ratio we need to find is
 \[\dfrac{\pi x^2}{\pi y^2} = \dfrac{x^2}{y^2}\]Draw the radii from the centers of the circles to $A$ and $B$. We can easily conclude that the $30^{\circ}$ belongs to the larger circle, and the $60$ degree arc belongs to the smaller circle. Therefore, $m\angle AO_1B = 30^{\circ}$ and $m\angle AO_2B = 60^{\circ}$. Note that $\Delta AO_2B$ is equilateral, so when chord AB is drawn, it has length $y$. Now, applying the Law of Cosines on $\Delta AO_1B$:
 \[y^2 = x^2 + x^2 - 2x^2\cos{30} = 2x^2 - x^2\sqrt{3} = (2 - \sqrt{3})x^2\]
 \[\dfrac{x^2}{y^2} = \dfrac{1}{2 - \sqrt{3}} = \dfrac{2 + \sqrt{3}}{4-3} = 2 + \sqrt{3}=\boxed{\textbf{(D)}}\](Solution by brandbest1)

 



\textbf{Solution 2}

Again, let the radius of the larger and smaller circles be $x$ and $y$, respectively, and let the centers of these circles be $O_1$ and $O_2$, respectively.  Let $X$ bisect segment $AB$.  Note that $\triangle AXO_1$ and $\triangle AXO_2$ are right triangles, with $\angle AO_1X=15^{\circ}$ and $\angle AO_2X=30^{\circ}$.  We have $\sin{15} = \dfrac{AX}{x}$ and $\sin{30} = \dfrac{AX}{y}$  and $\dfrac{x}{y} = \dfrac{\sin{30}}{\sin{15}}$. Since the ratio of the area of the larger circle to that of the smaller circle is simply $\dfrac{\pi x^2}{\pi y^2} = \left(\dfrac{x}{y} \right)^2 = \left(\dfrac{\sin{30}}{\sin{15}} \right)^2$, we just need to find $\sin{30}$ and $\sin{15}$. We know $\sin{30} = \dfrac{1}{2}$, and we can use the angle sum formula or half angle formula to compute $\sin{15} = \dfrac{\sqrt{6} - \sqrt{2}}{4}$. Plugging this into the previous expression, we get: 
 \[\left(\dfrac {x}{y} \right)^2 = \left(\dfrac{\dfrac{1}{2}}{\dfrac{\sqrt{6} - \sqrt{2}}{4}} \right)^2 = \left(\dfrac{\sqrt{6} + \sqrt{2}}{2} \right)^2 = 2 + \sqrt{3} =\boxed{\textbf{(D)}}\](Solution by kevin38017)

 



\textbf{Solution 3}

Let the radius of the smaller and larger circle be $r$ and $R$, respectively. We see that half the length of the chord is equal to $r \sin 30^{\circ}$, which is also equal to $R \sin 15^{\circ}$. Recall that $\sin 15^{\circ} = \frac{\sqrt{6} - \sqrt{2}}{4}$ and $\sin 30^{\circ} = \frac{1}{2}$. From this, we get $r = \frac{\sqrt{6} - \sqrt{2}}{2} R$, or $r^2 = \frac{8 - 2 \sqrt{12}}{4} R^2 = \left(2 - \sqrt{3}\right) R^2$, which is equivalent to $R^2 = \left(2 + \sqrt{3}\right) r^2$.

 (Solution by soy\_un\_chemisto)

 



\textbf{Solution 4}

As in the previous solutions let the radius of the smaller and larger circles be $r$ and $R$, respectively. Also, let their centers be $O_1$ and $O_2$, respectively. Now draw two congruent chords from points $A$ and $B$ to the end of the smaller circle, creating an isosceles triangle. Label that point $X$. Recalling the Inscribed Angle Theorem, we then see that $m\angle AXB = \frac{m\angle AO_1B}{2} = 30^{\circ}= m\angle AO_2B$. Based on this information, we can conclude that triangles $AXB$ and  $AO_2B$ are congruent via ASA Congruence.

 
Next draw the height of $AXB$ from $X$ to $AB$. Note we've just created a right triangle with hypotenuse $R$, base $\frac{r}{2}$, and height $\frac{r\sqrt{3}}{2} + r$Thus using the Pythagorean Theorem we can express $R^2$ in terms of $r$
 \[R^2 = (\frac{r}{2})^2 + (\frac{r\sqrt{3}}{2} + r)^2 = r^2 + \frac{r^2}{4} + \frac{3r^2}{4} + (2)(\frac{r\sqrt{3}}{2})(r) = 2r^2 + r^2\sqrt{3} = r^2(2 + \sqrt{3})\]

 We can now determine the ratio between the larger and smaller circles:
 \[\frac{Area [O_2]}{Area [O_1]} = \frac{\pi R^2}{\pi r^2} = \frac{\pi r^2(2 + \sqrt{3})}{\pi r^2} =\boxed{\textbf{(D) }  2 + \sqrt{3}}\]

 (Solution by derekxu)

 



\textbf{Solution 5}

Let the radius of the larger and smaller circles be $R$ and $r$, respectively, and let the centers of these circles be $O_1$ and $O_2$, respectively. Draw the radii and $AB$, and note that $AB=r$ because $\triangle ABO_1$ is equilateral. Also, $m\angle O_1AB = m\angle O_1BA = \frac{180^{\circ}-30^{\circ}}{2} = 75^{\circ}$. Then, mark point $X$ inside the larger circle such that $\angle AXB$ is a right angle and $AX = BX$. Notice that $\triangle ABX$ is a 45-45-90 triangle, so $AX = BX = \frac{r\sqrt{2}}{2}$. Now extend $BX$ to $AO_1$, and label the intersection $Y$. Since $m\angle YAB = 75^{\circ}$, $m\angle YAX = 75^{\circ} - 45^{\circ} = 30^{\circ}$ so this creates 30-60-90 triangle $\triangle AYX$. Therefore, $YX = \frac{r\sqrt{6}}{6}$ and $AY = \frac{r\sqrt{6}}{3}$. We can also see that $\angle YBO_1 = 75^{\circ}-45^{\circ} = 30^{\circ}=\angle YO_1B$, so $\triangle YO_1B$ is an isosceles triangle with $YB = YO_1$. So $YO_1 = \frac{r\sqrt{6}}{6} + \frac{r\sqrt{2}}{2}$. This means:
 \[R = AO_1 = \frac{r\sqrt{6}}{6} + \frac{r\sqrt{2}}{2} + \frac{r\sqrt{6}}{3} = \frac{r(\sqrt{6} + \sqrt{2})}{2}\]Then we can find the ratio:
 \[\frac{\pi R^2}{\pi r^2} = \frac{\pi [\frac{r(\sqrt{6} + \sqrt{2})}{2}]^2}{\pi r^2} = \frac{(\sqrt{6}+\sqrt{2})}{4}=\frac{8+4\sqrt{3}}{4}=\boxed{\textbf{(D) }  2 + \sqrt{3}}\]

 (Solution by weirdo)

 



\textbf{Solution 6 (using answer choices)}

We will estimate the answer using a wrong method then guess the correct answer choice.

 Let the radius of the larger and smaller circles be $R$ and $r$, respectively. Pretend line $AB$ is equal to the arc length of both circles, then 
 \[\frac{30^{\circ}}{360^{\circ}}\cdot 2\pi R=\frac{60^{\circ}}{360^{\circ}}\cdot 2\pi r\]
 \[R=2r\] and the answer to the problem is 
 \[\frac{\pi R^2}{\pi r^2} = 4\]But, this is only an estimate, the correct answer is not 4, but instead \textit{about} 4. We can asume that the closest other answer choice to 4 is correct: 
 \[\boxed{\textbf{(D) }2+\sqrt3}\]

 (Solution by FlareVa)

 


	\item (\textit{AMC 12 12A 2014}) A fancy bed and breakfast inn has $5$ rooms, each with a distinctive color-coded decor.  One day $5$ friends arrive to spend the night.  There are no other guests that night.  The friends can room in any combination they wish, but with no more than $2$ friends per room.  In how many ways can the innkeeper assign the guests to the rooms?


 $\textbf{(A) }2100\qquad \textbf{(B) }2220\qquad \textbf{(C) }3000\qquad \textbf{(D) }3120\qquad \textbf{(E) }3125\qquad$

 \textbf{Answer: }B



\textbf{Solution 1}

We can discern three cases.

 \textbf{Case 1:} Each room houses one guest. In this case, we have $5$ guests to choose for the first room, $4$ for the second, ..., for a total of $5!=120$ assignments.

 \textbf{Case 2:} Three rooms house one guest; one houses two. We have $\binom{5}{3}$ ways to choose the three rooms with $1$ guest, and $\binom{2}{1}$ to choose the remaining one with $2$. There are $5\cdot4\cdot3$ ways to place guests in the first three rooms, with the last two residing in the two-person room, for a total of $\binom{5}{3}\binom{2}{1}\cdot5\cdot4\cdot3=1200$ ways.

 \textbf{Case 3:} Two rooms house two guests; one houses one. We have $\binom{5}{2}$ to choose the two rooms with two people, and $\binom{3}{1}$ to choose one remaining room for one person. Then there are $5$ choices for the lonely person, and $\binom{4}{2}$ for the two in the first two-person room. The last two will stay in the other two-room, so there are $\binom{5}{2}\binom{3}{1}\cdot5\cdot\binom{4}{2}=900$ ways.

 In total, there are $120+1200+900=2220$ assignments, or $\boxed{\textbf{(B)}}$.

 (Solution by AwesomeToad16)

 



\textbf{Solution 2}

We can work in reverse by first determining the number of combinations in which there are more than $2$ friends in at least one room. There are three cases:

 \textbf{Case 1:} Three friends are in one room. Since there are $5$ possible rooms in which this can occur, we are choosing three friends from the five, and the other two friends can each be in any of the four remaining rooms,there are $5\cdot\binom{5}{3}\cdot4\cdot4 = 800$ possibilities.

 \textbf{Case 2:} Four friends are in one room. Again, there are $5$ possible rooms, we are choosing four of the five friends, and the other one can be in any of the other four rooms, so there are $5\cdot\binom{5}{4}\cdot4= 100$ possibilities.

 \textbf{Case 3:} Five friends are in one room. There are $5$ possible rooms in which this can occur, so there are $5$ possibilities.

 Since there are $5^5 = 3125$ possible combinations of the friends, the number fitting the given criteria is $3125 - (800+100+5) = \boxed{2220}$.

 


	\item (\textit{AMC 12 12A 2014}) Let $a<b<c$ be three integers such that $a,b,c$ is an arithmetic progression and $a,c,b$ is a geometric progression.  What is the smallest possible value of $c$?


 $\textbf{(A) }-2\qquad \textbf{(B) }1\qquad \textbf{(C) }2\qquad \textbf{(D) }4\qquad \textbf{(E) }6\qquad$

 \textbf{Answer: }C



\textbf{Solution 1}

We have $b-a=c-b$, so $a=2b-c$. Since $a,c,b$ is geometric, $c^2=ab=(2b-c)b \Rightarrow 2b^2-bc-c^2=(2b+c)(b-c)=0$. Since $a<b<c$, we can't have $b=c$ and thus $c=-2b$. Then our arithmetic progression is $4b,b,-2b$. Since $4b < b < -2b$, $b < 0$. The smallest possible value of $c=-2b$ is $(-2)(-1)=2$, or $\boxed{\textbf{(C)}}$.

 (Solution by AT)

 



\textbf{Solution 2}

Taking the definition of an arithmetic progression, there must be a common difference between the terms, giving us $(b-a) = (c-b)$. From this, we can obtain the expression $a = 2b-c$. Again, by taking the definition of a geometric progression, we can obtain the expression, $c=ar$ and $b=ar^2$, where r serves as a value for the ratio between two terms in the progression. By substituting $b$ and $c$ in the arithmetic progression expression with the obtained values from the geometric progression, we obtain the equation, $a=2ar^2-ar$ which can be simplified to $(r-1)(2r+1)=0$ giving us $r=1$ or $r=-1/2$. Thus, from the geometric progression, $a=a$, $b=1/4a$ and $c=-1/2a$. Looking at the initial conditions of $a<b<c$ we can see that the lowest integer value that would satisfy the above expressions is if $a = -4$, thus making $c=2$ or $\boxed{\textbf{(C)}}$ (Solution by thatuser)

 



\textbf{Solution 3}

By the definition of an arithmetic progression, we can label the terms $a$, $b$, and $c$, as $a$, $a+d$, and $a+2d$. Now, we have that $a$, $a+2d$, and $a+d$ form a geometric progression. Since a geometric ratio has a common ratio between terms, we have $(a+2d)/a = (a+d)/a+2d$. Cross multiplying, we obtain the equation $(a+2d)^2=a(a+d)$. Multiplying it out and cancelling terms, we are left with the quadratic equation $4d^2+3ad=0$. Solving for $d$ in terms of $a$, we get that $d=-3a/4$ or $d=0$. Looking at the problem, we know that the $d$ cannot be 0, therefore the arithmetic progression is $a, a/4, -a/2$, so we need to find the minimum value of $-a/2$ while $-a/2>a$. Looking at our progression, we realize that a must be a multiple of 4 so that every term is an integer. Substituting $a=-4$, since that would yield the smallest value of $-a/2$ which satisfies the conditions, we figure out that the answer is $\boxed{\textbf{(C)}}$.(Solution by i8Pie)

 



\textbf{Solution 4 (Directly solving for c)}

Using the set-up in the second solution, we can obtain $2b=a+c$ and $c^2=ab$. Directly solving for $c$, we obtain $c^2=a(\frac{a+c}{2})$ and thus $2c^2-ac-a^2=0$, a polynomial with respect to $c$. Solving using the quadratic formula yields results $c=a$ or $c=\frac{1}{2}a$. By definition, $a<b<c$ so $c=\frac{1}{2}a$. Solving for b yields $b=\frac{a}{4}$. The rest follows solution 2 and gives $c=2$ or $\boxed{\textbf{(C)}}$(Solution by cayleynotfermat1005)

 


	\item (\textit{AMC 12 12A 2014}) A five-digit palindrome is a positive integer with respective digits $abcba$, where $a$ is non-zero.  Let $S$ be the sum of all five-digit palindromes.  What is the sum of the digits of $S$?


 $\textbf{(A) }9\qquad \textbf{(B) }18\qquad \textbf{(C) }27\qquad \textbf{(D) }36\qquad \textbf{(E) }45\qquad$

 \textbf{Answer: }B



\textbf{Solution 1}

For each digit $a=1,2,\ldots,9$ there are $10\cdot10$ (ways of choosing $b$ and $c$) palindromes. So the $a$s contribute $(1+2+\cdots+9)(100)(10^4+1)$ to the sum.For each digit $b=0,1,2,\ldots,9$ there are $9\cdot10$ (since $a \neq 0$) palindromes. So the $b$s contribute $(0+1+2+\cdots+9)(90)(10^3+10)$ to the sum.Similarly, for each $c=0,1,2,\ldots,9$ there are $9\cdot10$ palindromes, so the $c$ contributes $(0+1+2+\cdots+9)(90)(10^2)$ to the sum.

 It just so happens that 
 \[(1+2+\cdots+9)(100)(10^4+1)+(1+2+\cdots+9)(90)(10^3+10)+(1+2+\cdots+9)(90)(10^2)=49500000\] so the sum of the digits of the sum is $\boxed{\textbf{(B)}\; 18}$.

 



\textbf{Solution 2}

Notice that $10001+ 99999 = 110000.$ In fact, ordering the palindromes in ascending order, we find that the sum of the nth palindrome and the nth to last palindrome is $110000.$ We have $9\cdot 10\cdot 10$ palindromes, or $450$ pairs of palindromes summing to $110000.$ Performing the multiplication gives $49500000$, so the sum $\boxed{\textbf{(B)}\; 18}$.

 



\textbf{Solution 3}

As shown above, there are a total of $900$ five-digit palindromes.  We can calculate their sum by finding the expected value of a randomly selected palindrome satisfying the conditions given, then multiplying it by $900$ to get our sum.  The expected value for the ten-thousands and the units digit is $\frac{1+2+3+\cdots+9}{9}=5$, and the expected value for the thousands, hundreds, and tens digit is $\frac{0+1+2+\cdots+9}{10}=4.5$.  Therefore our expected value is $5\times10^4+4.5\times10^3+4.5\times10^2+4.5\times10^1+5\times10^0=55,\!000$.  Since the question asks for the sum of the digits of the resulting sum, we do not need to keep the trailing zeros of either $55,\!000$ or $900$.  Thus we only need to calculate $55\times9=495$, and the desired sum is $\boxed{\textbf{(B) }18}$.

 



\textbf{Solution 4 (Variation of #2)}

First, allow $a$ to be zero, and then subtract by how much we overcount. We'll also sum each palindrome with its $\textit{complement}$. If $\overline{abcba}$ (the line means a, b, and c are digits and $abcba\ne a\cdot b\cdot c\cdot b\cdot a$) is a palindrome, then its complement is $\overline{defed}$ where $d=9-a$, $e=9-b$, $f=9-c$. Notice how every palindrome has a unique compliment, and that the sum of a palindrome and its complement is $99999$. Therefore, the sum of our palindromes is $99999\times (10^3/2)$. (There are $10^3/2$ pairs.)

 However, we have overcounted, as something like $05350$ $\textit{isn't}$ a palindrome by the problem's definition, but we've still included it. So we must subtract the sum of numbers in the form $\overline{0nmn0}$. By the same argument as before, these sum to $9990\times  (10^2/2)$. Therefore, the sum that the problem asks for is:

 
 \[500\times99999-50\times 9990\]
 \[=500\times99999-500\times 999\]
 \[=500(99999-999)\]
 \[=500\times 99000\]

 Since all we care about is the sum of the digits, we can drop the $0$'s.

 
 \[5\times99\]
 \[=5\times(100-1)\]
 \[=495\]

 And finally, $4+9+5=\boxed{\textbf{(B)}18}$



	\item (\textit{AMC 12 12A 2014}) The product $(8)(888\ldots 8)$, where the second factor has $k$ digits, is an integer whose digits have a sum of $1000$.  What is $k$?


 $\textbf{(A) }901\qquad \textbf{(B) }911\qquad \textbf{(C) }919\qquad \textbf{(D) }991\qquad \textbf{(E) }999\qquad$

 \textbf{Answer: }D



\textbf{Solution 1}

We can list the first few numbers in the form $8 \cdot (8....8)$

 (Hard problem to do without the multiplication, but you can see the pattern early on)

 $8 \cdot 8 = 64$
$8 \cdot 88 = 704$
$8 \cdot 888 = 7104$
$8 \cdot 8888 = 71104$
$8 \cdot 88888 = 711104$

 By now it's clear that the numbers will be in the form $7$, $k-2$ $1$'s, and $04$. We want to make the numbers sum to 1000, so $7+4+(k-2) = 1000$. Solving, we get $k = 991$, meaning the answer is $\boxed{\textbf{(D) } 991}$

 Another way to proceed is that we know the difference between the sum of the digits of each product and $k$ is always $9$, so we just do $1000-9=\boxed{\textbf{(D) } 991}$.

 \textbf{Proof of Solution 1} Since this solution won't fly on a proof-based competition, here's a proof that it's valid:

 We will call $x_k=8(888\dots8)$ with exactly $k$ $8$s. We then rewrite this more formally, as:

 
 \[x_k=8\biggr(\sum_{n=0}^{k-1}8(10)^n\biggr)\]
 \[=64\biggr(\sum_{n=0}^{k-1}(10)^n\biggr)\]
 \[=64\frac{10^{k}-1}{9}\]

 Then, finding a recursive formula, we get:

 
 \[x_{k+1}=64\times 10^{k}+x_k\]

 We will now use induction, Our base case will be $k=2$. It's easy to see that this becomes $x_2=704$. Then, the $k+1$ case: let $x_k=7111\dots104$ with $k-2$ $1$s. Then $x_{k+1}=64000\dots000+7111\dots104$. Adding these numbers, we get $x_{k+1}=71111\dots104$. 

 Summing these digits, we have $4+7+(k-2)=1000$, giving us $k=991$.

 



\textbf{Solution 2 (Educated Guessing if you have no time)}

We first note that $125 \cdot 8 = 1000$ and so we assume there are $125$ 8s.

 Then we note that it is asking for the second factor, so we subtract $1$(the original $8$ in the first factor).

 Now we have $125-1=124.$ The second factor is obviously a multiple of $124$.

 Listing the first few, we have $124, 248, 372, 496, 620, 744, 868, 992, 1116, 1240, ...$

 We notice that the 4th answer choice is 1 less than a 992(a multiple of 124.)

 Thus we make an educated guess that it is somehow less by 1, so we get $\fbox{(D)}$. ~mathboy282

 \textbf{Note (Must Read)} We were just lucky; this method is NOT reliable. Please note that this probably will not work for other problems; it is just a lucky coincidence.

 



\textbf{Solution 3}

You can write the expression as 8(8*10\textasciicircum{}k + 8*10\textasciicircum{}k-1...)

 This equals 64*10\textasciicircum{}k + 64*10\textasciicircum{}k-1...so if we add an example we see something.

 6400000000+

 0640000000+

 0064000000+

 0006400000+

 0000640000+

 0000064000+

 0000006400+

 0000000640+

 0000000064=

 7111111104

 You see that is just 7 - M 1s- 04

 There have to be 1000 - 7 - 4 = 989 1s.

 When you divide 7111... by 8, the answer, since 7 < 8, is 1 digit less.

 So we get 992 - 1 = 991, which is $\fbox{(D)}$. ~vedantinity

 


	\item (\textit{AMC 12 12A 2014}) A $4\times 4\times h$ rectangular box contains a sphere of radius $2$ and eight smaller spheres of radius $1$.  The smaller spheres are each tangent to three sides of the box, and the larger sphere is tangent to each of the smaller spheres.  What is $h$?


 
                    \begin{figure}[H]
                        \centering
                        \includegraphics[width=0.5\textwidth]{Problem/AMC_12_12A_2014_17_Problem_Figure_1.png}
                    \end{figure}
                    
$\textbf{(A) }2+2\sqrt 7\qquad \textbf{(B) }3+2\sqrt 5\qquad \textbf{(C) }4+2\sqrt 7\qquad \textbf{(D) }4\sqrt 5\qquad \textbf{(E) }4\sqrt 7\qquad$

 \textbf{Answer: }A



\textbf{Solution 1 (Light Coordinate Geometry)}

Let one of the corners be $(0, 0, 0)$. We can orient the box such that the center of the small sphere closest to the corner is $(1,1,1)$, and the center of the large sphere is $(2, 2, h/2)$. 

 Since the two spheres are tangent, the distance between their centers is $1+2 = 3$, so $\sqrt{(2-1)^2+(2-1)^2+(h/2-1)^2} = 3$. Solving, $h=2 + 2\sqrt{7}=\boxed{\textbf{(A)}}$




\textbf{Solution 2}

Let $A$ be the point in the same plane as the centers of the top spheres equidistant from said centers. Let $B$ be the analogous point for the bottom spheres, and let $C$ be the midpoint of $\overline{AB}$ and the center of the large sphere. Let $D$ and $E$ be the points at which line $AB$ intersects the top of the box and the bottom, respectively.

 Let $O$ be the center of any of the top spheres (you choose!). We have $AO=1\cdot\sqrt{2}$, and $CO=3$, so $AC=\sqrt{3^2-\left(\sqrt2\right)^2}=\sqrt{7}$. Similarly, $BC=\sqrt{7}$. $\overline{AD}$ and $\overline{BE}$ are clearly equal to the radius of the small spheres, $1$. Thus the total height is $AD+AC+BC+BE=2+2\sqrt7$, or $\boxed{\textbf{(A)}}$.

 



\textbf{Solution 3}


                    \begin{figure}[H]
                        \centering
                        \includegraphics[width=0.5\textwidth]{Solution/AMC_12_12A_2014_17_Solution_Figure_1.JPG}
                    \end{figure}
                    captionLet $A$ be the center of the large sphere and $C$ be the center of any small sphere. Let $D$ be a vertex of the rectangular prism closest to point $C$. Let $F$ be the point on the edge of the prism such that $\overline{DF}$ and $\overline{AF}$ are perpendicular. Let points $B$ and point $E$ lie on $\overline{AF}$ and $\overline{DF}$ respectively such that $\overline{CE}$ and $\overline{CB}$ are perpendicular at $C$. 

 
$AC$ is the radii of the spheres, so $AC=2+1=3$. $CE$ is the shortest length between the center of a small sphere and the edge of the prism, so $CE=\sqrt{2}$. Similarly, $AF=2\sqrt{2}$. Since $CEFB$ is a rectangle, $BF=CE=\sqrt{2}$. Since $AF=2\sqrt{2}$, $AB=AF - BF = \sqrt{2}$. Then, $BC=\sqrt{3^2-\left(\sqrt2\right)^2}=\sqrt{7}=EF$. $DE$ is the length from $C$ to the top of the prism or $1$. Thus, $DF=DE+EF=1+\sqrt{7}$. The prism is symmetrical, so $h=2DF=\boxed{\textbf{(A)}}$




\textbf{Solution 4}

Take a cross section and you will see that h is made up of the two radii of the circles plus some radical expression. The only choice satisfying this condition is $\boxed{\textbf{(A)}}$.

 


 


	\item (\textit{AMC 12 12A 2014}) The domain of the function $f(x)=\log_{\frac12}(\log_4(\log_{\frac14}(\log_{16}(\log_{\frac1{16}}x))))$ is an interval of length $\frac mn$, where $m$ and $n$ are relatively prime positive integers.  What is $m+n$?


 $\textbf{(A) }19\qquad \textbf{(B) }31\qquad \textbf{(C) }271\qquad \textbf{(D) }319\qquad \textbf{(E) }511\qquad$

 \textbf{Answer: }C



\textbf{Solution 1 (Generalization)}

For all real numbers $a,b,$ and $c$ such that $b>0$ and $b\neq1,$ note that:

 
            \begin{enumerate} 
        \item $\log_b a$ is defined if and only if $a>0.$

        \item For $0<b<1,$ we conclude that:
            \begin{itemize}<li>$\log_b a<c$ if and only if $a>b^c.$</li>
<li>$\log_b a>c$ if and only if $0<a<b^c.$</li>

            \end{itemize}
            For $b>1,$ we conclude that:

 
            \begin{itemize}<li>$\log_b a<c$ if and only if $0<a<b^c.$</li>
<li>$\log_b a>c$ if and only if $a>b^c.$</li>

            \end{itemize}
            
            \end{enumerate}
            Therefore, we have
 \begin{align*} \log_{\frac12}(\log_4(\log_{\frac14}(\log_{16}(\log_{\frac1{16}}x)))) \text{ is defined} &\implies \log_4(\log_{\frac14}(\log_{16}(\log_{\frac1{16}}x)))>0 \\ &\implies \log_{\frac14}(\log_{16}(\log_{\frac1{16}}x))>1 \\ &\implies 0<\log_{16}(\log_{\frac1{16}}x)<\frac14 \\ &\implies 1<\log_{\frac1{16}}x<2 \\ &\implies \frac{1}{256}<x<\frac{1}{16}. \end{align*}
The domain of $f(x)$ is an interval of length $\frac{1}{16}-\frac{1}{256}=\frac{15}{256},$ from which the answer is $15+256=\boxed{\textbf{(C) }271}.$
<u>\textbf{Remark}</u>

 This problem is quite similar to \href{https://artofproblemsolving.com/wiki/index.php?title=2004\_AMC\_12A\_Problems/Problem\_16}{2004 AMC 12A Problem 16}.

 



\textbf{Solution 2 (Substitution)}

For simplicity, let $a=\log_{\frac{1}{16}}{x},b=\log_{16}a,c=\log_{\frac{1}{4}}b$, and $d=\log_4c$.

 The domain of $\log_{\frac{1}{2}}x$ is $x \in (0, \infty)$, so $d \in (0, \infty)$.Thus, $\log_4{c} \in (0, \infty) \Rightarrow c \in (1, \infty)$.Since $c=\log_{\frac{1}{4}}b$ we have $b \in \left(0, \left(\frac{1}{4}\right)^1\right)=\left(0, \frac{1}{4}\right)$.Since $b=\log_{16}{a}$, we have $a \in (16^0,16^{1/4})=(1,2)$.Finally, since $a=\log_{\frac{1}{16}}{x}$, $x \in \left(\left(\frac{1}{16}\right)^2,\left(\frac{1}{16}\right)^1\right)=\left(\frac{1}{256},\frac{1}{16}\right)$.

 The length of the $x$ interval is $\frac{1}{16}-\frac{1}{256}=\frac{15}{256}$ and the answer is $\boxed{\textbf{(C) }271}$.

 



\textbf{Solution 3 (Calculus)}

The domain of $f(x)$ is the range of the inverse function $f^{-1}(x)=\left(\frac1{16}\right)^{16^{\left(\frac14\right)^{4^{\left(\frac12\right)^x}}}}$. Now $f^{-1}(x)$ can be seen to be strictly decreasing, since $\left(\frac12\right)^x$ is decreasing, so $4^{\left(\frac12\right)^x}$ is decreasing, so $\left(\frac14\right)^{4^{\left(\frac12\right)^x}}$ is increasing, so $16^{\left(\frac14\right)^{4^{\left(\frac12\right)^x}}}$ is increasing, therefore $\left(\frac1{16}\right)^{16^{\left(\frac14\right)^{4^{\left(\frac12\right)^x}}}}$ is decreasing.

 Therefore, the range of $f^{-1}(x)$ is the open interval $\left(\lim_{x\to\infty}f^{-1}(x), \lim_{x\to-\infty}f^{-1}(x)\right)$. We find:
 \begin{align*}    \lim_{x\to-\infty}\left(\frac1{16}\right)^{16^{\left(\frac14\right)^{4^{\left(\frac12\right)^x}}}}&= \lim_{a\to\infty}\left(\frac1{16}\right)^{16^{\left(\frac14\right)^{4^a}}}\\ &= \lim_{b\to\infty}\left(\frac1{16}\right)^{16^{\left(\frac14\right)^b}}\\ &= \left(\frac1{16}\right)^{16^0}\\ &= \frac{1}{16}. \end{align*}
Similarly,
 \begin{align*} \lim_{x\to\infty}\left(\frac1{16}\right)^{16^{\left(\frac14\right)^{4^{\left(\frac12\right)^x}}}}&=\left(\frac1{16}\right)^{16^{\left(\frac14\right)^{4^0}}}\\ &= \left(\frac1{16}\right)^{16^{\frac14}}\\ &= \left(\frac1{16}\right)^2\\ &= \frac{1}{256}. \end{align*}
Hence the range of $f^{-1}(x)$ (which is then the domain of $f(x)$) is $\left(\frac{1}{256},\frac{1}{16}\right)$ and the answer is $\boxed{\textbf{(C) }271}$.

 


	\item (\textit{AMC 12 12A 2014}) There are exactly $N$ distinct rational numbers $k$ such that $|k|<200$ and 
 \[5x^2+kx+12=0\] has at least one integer solution for $x$.  What is $N$?


 $\textbf{(A) }6\qquad \textbf{(B) }12\qquad \textbf{(C) }24\qquad \textbf{(D) }48\qquad \textbf{(E) }78\qquad$

 \textbf{Answer: }E



\textbf{Solution 1}

Factor the quadratic into
 \[\left(5x + \frac{12}{n}\right)\left(x + n\right) = 0\]where $-n$ is our integer solution. Then, 
 \[k = \frac{12}{n} + 5n,\]which takes rational values between $-200$ and $200$ when $|n| \leq 39$, excluding $n = 0$. This leads to an answer of $2 \cdot 39 = \boxed{\textbf{(E) } 78}$.

 



\textbf{Solution 2}

Solve for $k$ so 
 \[k=-\frac{12}{x}-5x.\]  Note that $x$ can be any integer in the range $[-39,0)\cup(0,39]$  so $k$ is rational with $\lvert k\rvert<200$.  Hence, there are $39+39=\boxed{\textbf{(E) } 78}.$




\textbf{Solution 3}

Plug in $k=200$ to find the upper limit. You will find the limit to be a number from $0<x<-1$ and one that is just below $-39.$ All the integer values from $-1$ to $-39$ can be attainable through some value of $k$. Since the question asks for the absolute value of $k$, we see that the answer is $39\cdot2 = \boxed{\textbf{(E)  }78.}$

 iron

 


	\item (\textit{AMC 12 12A 2014}) In $\triangle BAC$, $\angle BAC=40^\circ$, $AB=10$, and $AC=6$.  Points $D$ and $E$ lie on $\overline{AB}$ and $\overline{AC}$ respectively.  What is the minimum possible value of $BE+DE+CD$?


 $\textbf{(A) }6\sqrt 3+3\qquad \textbf{(B) }\dfrac{27}2\qquad \textbf{(C) }8\sqrt 3\qquad \textbf{(D) }14\qquad \textbf{(E) }3\sqrt 3+9\qquad$

 \textbf{Answer: }D



\textbf{Solution 1}

Let $C_1$ be the reflection of $C$ across $\overline{AB}$, and let $C_2$ be the reflection of $C_1$ across $\overline{AC}$.  Then it is well-known that the quantity $BE+DE+CD$ is minimized when it is equal to $C_2B$.  (Proving this is a simple application of the triangle inequality; for an example of a simpler case, see Heron's Shortest Path Problem.)  As $A$ lies on both $AB$ and $AC$, we have $C_2A=C_1A=CA=6$.  Furthermore, $\angle CAC_1=2\angle CAB=80^\circ$ by the nature of the reflection, so $\angle C_2AB=\angle C_2AC+\angle CAB=80^\circ+40^\circ=120^\circ$.  Therefore by the Law of Cosines 
 \[BC_2^2=6^2+10^2-2\cdot 6\cdot 10\cos 120^\circ=196\implies BC_2=\boxed{\textbf{(D) }14}.\]




\textbf{Solution 2}

In $\triangle BAC$, the three lines look like the Chinese character 又. Let $\triangle DEA$, $\triangle CDA$, and $\triangle BEA$ have bases $DE$, $CD$, and $BE$ respectively. Then, $\triangle DEA$ has the same side $DA$ as $\triangle CDA$ and the same side $EA$ as $\triangle BEA$. Connect all three triangles with $\triangle DEA$ in the center and the two triangles sharing one of its sides. Then, $\pentagon BACDE$ is formed with $BE+DE+CD$ forming the base. 

 Intuitively, the pentagon's base is minimized when all three bottom sides are collinear. This is simply the original $\triangle BAC$ except that $\angle BAC =120^\circ$. (In $\triangle DEA$, $\triangle CDA$, and $\triangle BEA$, $\angle A = 40^\circ$, and the three triangles connect at $A$ to form the pentagon). Thus, $m\angle BAC = 40 \cdot 3$).

 
$BC$ in this new triangle is then the minimum of $BE+DE+CD$. Applying law of cosines, $BC=\sqrt{6^2+10^2-2(6)(10)\cos (120^\circ)}=\sqrt{196}=14 \implies \boxed{\textbf{(D) }14}$
\begin{verbatim}
Would prime notation be clearer?
\end{verbatim}





\textbf{Solution 3}


                    \begin{figure}[H]
                        \centering
                        \includegraphics[width=0.5\textwidth]{Solution/AMC_12_12A_2014_20_Solution_Figure_1.png}
                    \end{figure}
                    

 (Diagram by shihan)Reflect $C$ across $AB$ to $C'$. Similarly, reflect $B$ across $AC$ to $B'$. Clearly, $BE = B'E$ and $CD = C'D$. Thus, the sum $BE + DE + CD = B'E + DE + C'D$. This value is minimized when $B'$, $C'$, $D$ and $E$ are collinear. To finish, we use the law of cosines on the triangle $AB'C'$: $B'C' = \sqrt{6^2 + 10^2 - 2(6)(10)\cos 120^{\circ}} = \boxed{\textbf{(D) } 14}.$

 


 Note (note by JiYang): $\quad \because \angle BAC \textless \frac{\pi}{3}$ $\qquad \therefore \angle C'AB' \textless \pi$ $\qquad \therefore$ the line C'B' would pass through $\Delta ABC$

 \textbf{Remark} This problem is similar to \href{https://en.wikipedia.org/wiki/Fagnano\%27s\_problem}{Fagnano's Problem}.

 


	\item (\textit{AMC 12 12A 2014}) For every real number $x$, let $\lfloor x\rfloor$ denote the greatest integer not exceeding $x$, and let 
 \[f(x)=\lfloor x\rfloor(2014^{x-\lfloor x\rfloor}-1).\]  The set of all numbers $x$ such that $1\leq x<2014$ and $f(x)\leq 1$ is a union of disjoint intervals.  What is the sum of the lengths of those intervals?


 $\textbf{(A) }1\qquad \textbf{(B) }\dfrac{\log 2015}{\log 2014}\qquad \textbf{(C) }\dfrac{\log 2014}{\log 2013}\qquad \textbf{(D) }\dfrac{2014}{2013}\qquad \textbf{(E) }2014^{\frac1{2014}}\qquad$

 \textbf{Answer: }A



\textbf{Solution}

Let $\lfloor x\rfloor=k$ for some integer $1\leq k\leq 2013$.  Then we can rewrite $f(x)$ as $k(2014^{x-k}-1)$.  In order for this to be less than or equal to $1$, we need $2014^{x-k}-1\leq\dfrac1k\implies x\leq k+\log_{2014}\left(\dfrac{k+1}k\right)$.  Combining this with the fact that $\lfloor x\rfloor =k$ gives that $x\in\left[k,k+\log_{2014}\left(\dfrac{k+1}k\right)\right]$, and so the length of the interval is $\log_{2014}\left(\dfrac{k+1}k\right)$.  We want the sum of all possible intervals such that the inequality holds true; since all of these intervals must be disjoint, we can sum from $k=1$ to $k=2013$ to get that the desired sum is 
 \[\sum_{i=1}^{2013}\log_{2014}\left(\dfrac{i+1}i\right)=\log_{2014}\left(\prod_{i=1}^{2013}\dfrac{i+1}i\right)=\log_{2014}\left(\dfrac{2014}1\right)=\boxed{1\textbf{ (A)}}.\]



	\item (\textit{AMC 12 12A 2014}) The number $5^{867}$ is between $2^{2013}$ and $2^{2014}$.  How many pairs of integers $(m,n)$ are there such that $1\leq m\leq 2012$ and 
 \[5^n<2^m<2^{m+2}<5^{n+1}?\]
$\textbf{(A) }278\qquad \textbf{(B) }279\qquad \textbf{(C) }280\qquad \textbf{(D) }281\qquad \textbf{(E) }282\qquad$

 \textbf{Answer: }B



\textbf{Solution 1}

Between any two consecutive powers of $5$ there are either $2$ or $3$ powers of $2$ (because $2^2<5^1<2^3$). Consider the intervals $(5^0,5^1),(5^1,5^2),\dots (5^{866},5^{867})$. We want the number of intervals with $3$ powers of $2$.

 From the given that $2^{2013}<5^{867}<2^{2014}$, we know that these $867$ intervals together have $2013$ powers of $2$. Let $x$ of them have $2$ powers of $2$ and $y$ of them have $3$ powers of $2$. Thus we have the system
 \[x+y=867\]
 \[2x+3y=2013\]from which we get $y=279$, so the answer is $\boxed{\textbf{(B)}}$.

 



\textbf{Solution 2}

The problem is asking for between how many consecutive powers of $5$ are there $3$ power of $2$s

 There can be either $2$ or $3$ powers of $2$ between any two consecutive powers of $5$, $5^n$ and $5^{n+1}$. 

 The first power of $2$ is between $5^n$ and $2 \cdot 5^n$.

 The second power of $2$ is between $2 \cdot 5^n$ and $4 \cdot 5^n$.

 The third power of $2$ is between $4 \cdot 5^n$ and $8 \cdot 5^n$, meaning that it can be between $5^n$ and $5^{n+1}$ or not.

 If there are only $2$ power of $2$s between every consecutive powers of $5$ up to $5^{867}$, there would be $867\cdot 2 = 1734$ power of $2$s. However, there are $2013$ powers of $2$ before  $5^{867}$, meaning the answer is $2013 - 1734 = \boxed{\textbf{(B)}279}$.

 


	\item (\textit{AMC 12 12A 2014}) The fraction 
 \[\dfrac1{99^2}=0.\overline{b_{n-1}b_{n-2}\ldots b_2b_1b_0},\] where $n$ is the length of the period of the repeating decimal expansion.  What is the sum $b_0+b_1+\cdots+b_{n-1}$?


 $\textbf{(A) }874\qquad \textbf{(B) }883\qquad \textbf{(C) }887\qquad \textbf{(D) }891\qquad \textbf{(E) }892\qquad$

 \textbf{Answer: }B



\textbf{Solution 1}

the fraction $\dfrac{1}{99}$ can be written as 
 \[\sum^{\infty}_{n=1}\dfrac{1}{10^{2n}}\].similarly the fraction $\dfrac{1}{99^2}$ can be written as $\sum^{\infty}_{m=1}\dfrac{1}{10^{2m}}\sum^{\infty}_{n=1}\dfrac{1}{10^{2n}}$ which is equivalent to 
 \[\sum^{\infty}_{m=1}\sum^{\infty}_{n=1} \dfrac{1}{10^{2(m+n)}}\]and we can see that for each $n+m=k$ there are $k-1$ $(n,m)$ combinations so the above sum is equivalent to:
 \[\sum^{\infty}_{k=2}\dfrac{k-1}{10^{2k}}\]we note that the sequence starts repeating at $k = 102$yet consider 
 \[\sum^{101}_{k=99}\dfrac{k-1}{10^{2k}}=\dfrac{98}{{10^{198}}}+\dfrac{99}{{10^{200}}}+\dfrac{100}{10^{{202}}}=\dfrac{1}{10^{198}}(98+\dfrac{99}{100}+\dfrac{100}{10000})=\dfrac{1}{10^{198}}(98+\dfrac{99}{100}+\dfrac{1}{100})=\dfrac{1}{10^{198}}(98+\dfrac{100}{100})=\dfrac{1}{10^{198}}(99)\]so the decimal will go from 1 to 99 skipping the number 98and we can easily compute the sum of the digits from 0 to 99 to be 
 \[45\cdot10\cdot2=900\] subtracting the sum of the digits of 98 which is 17 we get
 \[900-17=883\textbf{(B) }\qquad\]




\textbf{Solution 2}

$\frac{1}{99^2}\\\\ =\frac{1}{99} \cdot \frac{1}{99}\\\\ =\frac{0.\overline{01}}{99}\\\\ =0.\overline{00010203...9799}$

 So, the answer is $0+0+0+1+0+2+0+3+...+9+7+9+9=2\cdot10\cdot\frac{9\cdot10}{2}-(9+8)$ or $\boxed{\textbf{(B)}\ 883}$.

 There are two things to notice here. First, $\frac{1}{99}$ has a very simple and unique decimal expansion, as shown. Second, for $\frac{0.\overline{01}}{99}$ to itself produce a repeating decimal, $99$ has to evenly divide a sufficiently extended number of the form $101010101..$. This number will have $99$ ones (197 digits in total), as to be divisible by $9$ and $11$. The enormity of this number forces us to look for a pattern, and so we divide out as shown. Indeed, upon division (seeing how the remainders always end in "501" or "601" or, at last, "9801"), we find the repeating part $.00010203040506070809101112131415....9799$. If we wanted to further check our pattern, we could count the total number of digits in our quotient (not counting the first three): 195. Since $99<100$, multiplying by it will produce either $1$ or $2$ extra digits, so our quotient passes the test.

 Or note that the 1 has to be carried when you get to 100 so the 99 becomes 00 and the 1 gets carried again to 98 which becomes 99. 

 



\textbf{Solution 3 (kinda bad)}


 \[99^2=9801\] Does 9801 look familiar? The fraction $\frac{1}{9801}$ is a very famous fraction, because when expressed as a decimal, it equals 
 \[0.0001020304050607080910...9091929394959697990001020304...\] (the numbers from $00-99 \textbf{ EXCEPT } 98$). You should look at the above solutions for the reason. The sum of the digits of $98$ is $17$, so therefore, adding 17 to all of the choices should be a nice number. $\text{(B)}$ becomes $883+17=900$, which is very nice, so our answer should be $\boxed{\textbf{(B)}\ 883}$.

 


	\item (\textit{AMC 12 12A 2014}) Let $f_0(x)=x+|x-100|-|x+100|$, and for $n\geq 1$, let $f_n(x)=|f_{n-1}(x)|-1$.  For how many values of $x$ is $f_{100}(x)=0$?


 $\textbf{(A) }299\qquad \textbf{(B) }300\qquad \textbf{(C) }301\qquad \textbf{(D) }302\qquad \textbf{(E) }303\qquad$

 \textbf{Answer: }C



\textbf{Solution 1}

1. Draw the graph of $f_0(x)$ by dividing the domain into three parts.
                    \begin{figure}[H]
                        \centering
                        \includegraphics[width=0.5\textwidth]{Solution/AMC_12_12A_2014_24_Solution_Figure_1.png}
                    \end{figure}
                     

 2. Apply the recursive rule a few times to find the pattern.

 \textbf{Note:} $f_n(x) = |f_{n-1}(x)| - 10$ is used to enlarge the difference, but the reasoning is the same.
                    \begin{figure}[H]
                        \centering
                        \includegraphics[width=0.5\textwidth]{Solution/AMC_12_12A_2014_24_Solution_Figure_2.png}
                    \end{figure}
                    

 3. Extrapolate to $f_{100}$. Notice that the summits start $100$ away from $0$ and get $1$ closer each iteration, so they reach $0$ exactly at $f_{100}$.
                    \begin{figure}[H]
                        \centering
                        \includegraphics[width=0.5\textwidth]{Solution/AMC_12_12A_2014_24_Solution_Figure_3.png}
                    \end{figure}
                    
$f_{100}(x)$ reaches $0$ at $x = -300$, then zigzags between $0$ and $-1$, hitting $0$ at every even $x$, before leaving $0$ at $x = 300$.

 This means that $f_{100}(x) = 0$ at all even $x$ where $-300 \le x \le 300$. This is a $601$-integer odd-size range with even numbers at the endpoints, so just over half of the integers are even, or $\frac{601+1}{2} = \boxed{\textbf{(C) }301}$.(Revised by Flamedragon & Jason,C & \href{https://artofproblemsolving.com/wiki/index.php?title=User:Emerald\_block}{emerald\_block})

 



\textbf{Solution 2}

First, notice that the recursive rule moves the current value $1$ closer to $0$. Upon reaching $0$, it alternates between $-1$ and $0$. This means that $f_{100}(x) = 0$ exactly when $|f_0(x)| \le 100$ (to reach $0$ in time) and $f_0(x)$ is even (so $f_{100}(x) \ne -1$).

 Casework each part of $f_0(x)$ (where the expressions in the absolute values do not change sign):
 \[x \le -100 \implies f_0(x) = x-(x-100)+(x+100) = x+200\]so even $-300 \le x \le -100$ work.
 \[-100 \le x \le 100 \implies f_0(x) = x-(x-100)-(x+100) = -x\]so even $-100 \le x \le 100$ work.
 \[100 \le x \implies f_0(x) = x+(x-100)-(x+100) = x-200\]so even $100 \le x \le 300$ work.

 Putting these together, all even $x$ where $-300 \le x \le 300$ work. So, the answer is $2\cdot150+1 = \boxed{\textbf{(C)}\ 301}$.~revised by \href{https://artofproblemsolving.com/wiki/index.php?title=User:Emerald\_block}{emerald\_block}

 



\textbf{Solution 3}

Note $f_{100}(x) = 0$ when $|f_{99}(x)| -1$ = 0. This occurs when $f_{99}(x) = \pm 1$. 

 Then, repeating this process, we note $f_{99}(x) = \pm 1 \implies |f_{98}(x)| = 0, 2$, and hence $f_{98}(x) = 0, \pm 2$.

 Similarly, $f_{97}(x) = \pm 1, \pm 3$. Extrapolating this pattern, we must have $f_{0}(x) = 0, \pm 2$, $\dots$, $\pm 100$. Then, drawing the graph of $f_0$, 
                    \begin{figure}[H]
                        \centering
                        \includegraphics[width=0.5\textwidth]{Solution/AMC_12_12A_2014_24_Solution_Figure_4.png}
                    \end{figure}
                     we note for each of $0, \pm 2$, $\dots$, $\pm 98$, there are three solutions. For $\pm 100$, there is exactly $2$ solutions.

 So, the total amount of solutions is $99 \cdot 3 + 2 \cdot 2 = \boxed{\textbf{(C) }301}$



	\item (\textit{AMC 12 12A 2014}) The parabola $P$ has focus $(0,0)$ and goes through the points $(4,3)$ and $(-4,-3)$.  For how many points $(x,y)\in P$ with integer coordinates is it true that $|4x+3y|\leq 1000$?


 $\textbf{(A) }38\qquad \textbf{(B) }40\qquad \textbf{(C) }42\qquad \textbf{(D) }44\qquad \textbf{(E) }46\qquad$

 \textbf{Answer: }B



\textbf{Solution}

The parabola is symmetric through $y=- \frac{4}{3}x$, and the common distance is $5$, so the directrix is the line through $(1,7)$ and $(-7,1)$, which is the line 
 \[3x-4y = -25.\] Using the point-line distance formula, the parabola is the locus 
 \[x^2+y^2 = \frac{\left\lvert 3x-4y+25 \right\rvert^2}{3^2+4^2}\] which rearranges to $(4x+3y)^2 = 25(6x-8y+25)$.

 Let $m = 4x+3y \in \mathbb Z$, $\left\lvert m \right\rvert \le 1000$.  Put $m = 25k$ to obtain
 \[25k^2 = 6x-8y+25\]
 \[25k = 4x+3y.\]and accordingly we find by solving the system that $x = \frac{1}{2} (3k^2-3) + 4k$ and $y = -2k^2+3k+2$.

 One can show that the values of $k$ that make $(x,y)$ an integer pair are precisely odd integers $k$.*  For $\left\lvert 25k \right\rvert \le 1000$ this is $k= -39,-37,-35,\dots,39$, so $40$ values work and the answer is $\boxed{\textbf{(B)}}$.

 (Solution by C-273)

 
            \begin{itemize}
        \item You can do this by noting 2 has to divide $(3k^2-3)$, so $3k^2$ is congruent to 3 modulo 2. Dividing by 3 (since 2 and 3 are coprime), $k^2$ is odd, so $k$ must be odd. - (Orion 2010)
            \end{itemize}
            



\textbf{Solution 2}

The axis of $P$ is inclined at an angle $\theta$ relative to the coordinate axis, where $\tan\theta = \frac 34$. We rotate the coordinate axis by angle $\theta$ anti-clockwise, so that the parabola now has a vertical symmetry axis relative to the rotated coordinates. Let $(\widetilde{x}, \widetilde{y})$ be the coordinates in the rotated system. Then $(x,y)$ and $(\widetilde{x}, \widetilde{y})$ are related by 
 \begin{align}   \nonumber x = \widetilde{x}\cos\theta -\widetilde{y}\sin\theta &= \frac 45 \widetilde{x} - \frac 35 \widetilde{y}, \\    y = \widetilde{x}\sin\theta +\widetilde{y}\cos\theta &= \frac 35 \widetilde{x} + \frac 45 \widetilde{y} \end{align}In the rotated coordinate system, the parabola has focus at $(0,0)$ and the two points on it are at $(5,0)$ and $(-5,0)$. Therefore, the directrix is $\widetilde{y}=\pm 5$; we can, WLOG, choose $\widetilde{y}=-5$. For a point on the parabola, it is equidistant from the focus and directrix, so the equation of the parabola is 
 \begin{align}\tag{2}   \widetilde{x}^2+\widetilde{y}^2 = (\widetilde{y}+5)^2 \qquad  &\Longrightarrow\qquad \widetilde{y} = \frac{1}{10}(\widetilde{x}^2-25) \end{align}From $(1)$ we have $|4x+3y|=5\widetilde{x}$, so we need $|\widetilde{x}|<200$. Substituting $(2)$ in $(1)$, we get
 \begin{align*}   50x &= 40 \widetilde{x} - 3 \widetilde{x}^2 + 75, \\   50y &= 30 \widetilde{x} + 4 \widetilde{x}^2 - 100 \end{align*}
For $x$ to be an integer $\widetilde{x}$ must be a multiple of 5; setting $\widetilde{x}=5a$ we get 
 \[2x = 8a - 3 a^2 + 3\]Now we need $a$ to be odd, i.e. $\widetilde{x}=5a$ is an odd multiple of $5$, in which case we get $y = 3 a + 2 a^2 - 2$, which is also an integer. The values that satisfy the given conditions correspond to $\widetilde{x}= \{\pm 5\cdot (2k-1)\mid k = 0, 1, \ldots , 19 \}={-195, -185, -175, ..., 195}$, and there are $\boxed{\textbf{(B)} \ 40}$ such numbers.

 (Solution by Shaddoll)

 


\end{enumerate}